\section{Discussion: design tradeoffs}
\label{sec:discussion}

This section collects, in one place, the design choices that
shaped the library, paired with what each gained and what each
cost.  We have referred to these tradeoffs locally throughout the
chapter; the goal here is to make the global picture explicit so
that a downstream user can decide whether the library's contract
matches their needs.

\subsection{Finite, discrete probability}
\label{sec:discussion:finite}

\paragraph{Choice.}\ Every distribution is a \TT{PMF}; every
strategy/outcome carrier is a \TT{Fintype} or carries a
\TT{Fintype} hypothesis on the relevant theorem.  No
$\sigma$-algebra is ever invoked.

\paragraph{Gained.}\ Every strategic-form definition is a
$\sum$-of-products identity that the library's \TT{simp} set can
manipulate.  No measurability hypotheses propagate through
strategy spaces, profile distributions, or outcome distributions.
Equational reasoning under the \TT{PMF} monad is mechanically
straightforward.  Compilation between languages is a calculation
on finite sums and Dirac/Bind compositions.

\paragraph{Lost.}\ The cost of mechanizing measure theory and
the catalogue of classical results that fall outside the discrete
contract are described in detail in
\S\ref{sec:prelim:pmf}; we do not repeat them here.  What §8 adds
is the consequence for downstream consumers:\ a protocol or
auction layer that builds on this library inherits the discrete
contract and must either rephrase its objects in finite/discrete
shape (e.g.\ a discretized bid grid, a finite type lattice) or
build a parallel measure-theoretic infrastructure on top of
mathlib and accept the corresponding proof obligations.  In
practice this is a hard contract:\ existing downstream users of
the library work in the discrete idiom uniformly, and the
library's design does not provide an automatic bridge to a
continuous version.  An honest reading of \S\ref{sec:prelim:pmf}
should leave the consumer with a sense of which classical
theorems they can call into and which they cannot.

\subsection{Protocol/preference split}
\label{sec:discussion:protocol-preference}

\paragraph{Choice.}\ \TT{GameForm} is utility-free; \TT{KernelGame}
adds utility.  Solution concepts exist in two forms:\ EU-specific
predicates on \TT{KernelGame} and preference-parameterized
predicates (\TT{*For}) on \TT{GameForm}.

\paragraph{Gained.}\ Theorems about strategic structure (Kuhn,
information-set machinery, the deviation hierarchy) state and
prove once at the \TT{GameForm} layer, parametric in any
preference preorder.  EU is one preference preorder; lexicographic,
maximin, or non-EU preferences are other preorders that get the
same theorems.  The library is therefore reusable for, e.g.,
ordinal-preference solution concepts without rewriting the
strategic skeleton.

\paragraph{Lost.}\ One additional indirection for \emph{generic}
users.  A user who only ever cares about expected utility writes
\TT{G.IsNash $\sigma$} and is done; a user who wants to state a
generic theorem parametric in the preference must write
\TT{G.IsNashFor pref $\sigma$} and supply the preference
explicitly.  Two versions of every solution concept exist
(\TT{IsNash}, \TT{IsNashFor}, etc.), connected by \TT{\_iff\_}
bridges; the \TT{simp} set has to know about both.  In practice the
EU-only ergonomics dominates daily use and the generic version is
visible only to library authors stating preference-parametric
theorems.

\paragraph{Why.}\ The split is motivated by classical
mechanism-design practice~\citep{Myerson1991}, where the
``social-choice rule + payment scheme = mechanism'' decomposition
is exactly the protocol/preference split made structural.  Once we
saw that several classical theorems (dominance $\Rightarrow$ Nash;
correlated $\Rightarrow$ coarse correlated; the deviation hierarchy)
were preference-agnostic, building them at the \TT{GameForm} layer
became the natural design.

\subsection{The \texttt{*For} style versus typeclass overloading}
\label{sec:discussion:for-style}

\paragraph{Choice.}\ Preferences are passed as explicit
\TT{Prop}-valued arguments to \TT{*For} solution concepts; they
are not registered as typeclass instances on \TT{KernelGame}.

\paragraph{Gained.}\ Several library theorems quantify over two
preferences at once (monotonicity of $\IsNashFor$ in the
preference; Pareto dominance, which takes a $(\Pref, \sPref)$
pair; strict-implies-weak bridges; the
transport theorems of \S\ref{sec:semantic-core:transport}, which
relate a preference either to its pullback along an outcome
bijection or to its projection through a common observed outcome
view).
With explicit-argument preferences each of these is one line; with
a typeclass-based design each would force ad-hoc local instances
and \TT{let}-binding tricks at the call site.

\paragraph{Lost.}\ Slight verbosity at the call site for the
\emph{generic} user:\ they write \TT{G.IsNashFor G.euPref}~\Profile\
instead of \TT{G.IsNash}~\Profile\ whenever they instantiate the
generic version at EU.  The EU-only user is unaffected — the
EU-specific predicates (\TT{IsNash}, \TT{IsBestResponse},
\TT{IsCorrelatedEq}, …) are exposed directly and the
\TT{\_iff\_} bridges make round-tripping free for the EU case.

\paragraph{Why.}\ The honest version of the argument is not that
typeclasses are categorically wrong here — a \TT{[Preference G]}
typeclass with a default \TT{euPref} instance would in fact make
the EU-only case ergonomically uniform.  Rather, the multi-preference
theorems listed above fit explicit-argument predicates naturally
and fit typeclasses awkwardly, and the EU-only ergonomic is
recovered cheaply by the EU-specific layer.  The library does
use a typeclass where one is the right fit:\ \TT{PrefPreorder
pref} asserts that a preference is reflexive and transitive — a
property of the relation, not a choice between relations, and
Lean's inference resolves it without diamond/multi-preference
issues.

\subsection{Solution concepts on \texttt{KernelGame} only}
\label{sec:discussion:kg-only}

\paragraph{Choice.}\ \TT{IsNash}, \TT{IsDominant}, \TT{IsBR},
\TT{IsCorrelatedEq}, etc.\ are defined on \TT{KernelGame} and on
\TT{GameForm}.  They are \emph{not} re-defined on each
representation (NFG, EFG, MAID, MultiRound, FOSG); concrete
representations bridge to the central definitions via $\toKG$.

\paragraph{Gained.}\ The library has one canonical
``\texttt{IsNash} of an EFG game'' — the \TT{IsNash} of its
compiled kernel game.  Theorems about Nash on EFGs do not have to
be re-proved each time a new EFG-like representation is added;
they are corollaries of the core theorems via composition with
$\compileTo$.  This is the workhorse of the bridge architecture
of \S\ref{sec:languages:bridges}.

\paragraph{Lost.}\ Concrete representations occasionally have a
representation-specific solution concept that is more natural in
their syntax than in the kernel-game image (subgame-perfect
equilibrium for EFGs, perfect Bayesian equilibrium for Bayesian
games).  These exist as separate definitions (e.g.\
\TT{EFG.IsSubgamePerfectEq}); the library proves their relation
to the central \TT{IsNash} as a theorem rather than building it
into the definition.

\subsection{Noncomputable definitions and classical reasoning}
\label{sec:discussion:noncomputable}

\paragraph{Choice.}\ Many definitions are \TT{noncomputable},
particularly those that invoke real-number operations, the
classical disjunction \TT{Function.update} on abstract index types,
or fixed-point theorems.  We do not attempt to expose computable
witnesses for, e.g., \TT{mixed\_nash\_exists}.

\paragraph{Gained.}\ Proofs are shorter and more uniform.  Classical
reasoning is locally invoked via \TT{open Classical in} and is
flagged at every use.  Real-number arithmetic uses mathlib's full
real-number API without restriction.

\paragraph{Lost.}\ The library is not a Nash-equilibrium algorithm.
A \TT{\#eval}-style execution will not produce a witness from
\TT{mixed\_nash\_exists}.  Constructive consumers must write their
own algorithm and prove it converges to a Nash equilibrium of the
same kernel game.

\paragraph{Why.}\ The library targets formalization of theorems,
not algorithm certification.  Brouwer's fixed-point theorem is
inherently non-constructive; replacing it with a constructive
substitute (Sperner's lemma) would change the proof shape of
\Cref{thm:nash-mixed} without changing the statement, and would
not improve the user's experience unless they happened to be
running the Lean kernel as an evaluator.  We treat constructive
existence as a separate, downstream concern.

\subsection{Distributional, not deterministic, EFG semantics}
\label{sec:discussion:efg-distributional}

\paragraph{Choice.}\ \TT{GameTree.evalDist} returns a \TT{PMF}
$\Outcome$ given a \emph{behavioral} profile.  No deterministic
``play'' function exists at the EFG level; pure strategies are
lifted to behavioral via point-mass distributions.

\paragraph{Gained.}\ Chance nodes integrate without ad-hoc
adapters.  Kuhn's theorem becomes a structural identity
$\evalDist(\beta) = (\mathrm{realize}(\mu)) \mathbin{\bind}
\evalDist^{\mathrm{pure}}$ rather than a probabilistic
equivalence, and the full game-form $\toGF$ has the
$\K : \prod_i \Strat_i \to \PMF{\Outcome}$ shape natively.

\paragraph{Lost.}\ Reasoning about a single play of an EFG (a
single deterministic execution path) requires unfolding a
\TT{PMF}.  This is a wash:\ traditional textbook reasoning quietly
assumes determinism while writing distributional formulas, so
making the distributional semantics primary forces explicit
bookkeeping in places where the textbook is implicit.

\subsection{\texorpdfstring{$\mathrm{Fin}\,n$}{Fin n} player types in bridges}
\label{sec:discussion:fin-n}

\paragraph{Choice.}\ Bridge theorems often canonicalize the
player type to $\mathrm{Fin}\,n$ for the appropriate $n$.

\paragraph{Gained.}\ A canonical concrete carrier on which to
state the bridge.  Two bridges into and out of $\mathrm{Fin}\,n$
compose without re-indexing in the middle.

\paragraph{Lost.}\ Source-game player types (e.g.\
$\mathrm{Bool}$ for two-player games) must be transported across
an equivalence \TT{Bool $\simeq$ Fin 2} to use the bridge.

\paragraph{Why.}\ \TT{KernelGame.reindex}
(\S\ref{sec:semantic-core:reindex}) makes the transport a
one-line operation.  Treating $\mathrm{Fin}\,n$ as the canonical
finite-cardinality witness lets every bridge avoid having to
pick one for itself.

\subsection{What we did not formalize and why}
\label{sec:discussion:scope}

For honesty's sake we list what is missing.

\begin{itemize}\itemsep2pt
\item \emph{Repeated and stochastic games at infinite horizon.}
The library's repeated games are bounded.  The folk-theorem-style
results (Friedman, Fudenberg–Maskin) require infinite or
discounted horizons, which interact with measure theory and
which we have not formalized.

\item \emph{Continuous strategy / type / outcome spaces.}  See
\S\ref{sec:discussion:finite}.

\item \emph{Computational complexity.}  Nash existence is
formalized; PPAD-completeness~\citep{DaskalakisGoldbergPapadimitriou2009}
is not.  Approximate-equilibrium algorithmic results are not
formalized.

\item \emph{Cooperative game theory beyond Shapley values.}  The
library defines the Shapley value and proves null-player,
additivity, linearity, and symmetry lemmas.  The full uniqueness
axiomatization, the core, the nucleolus, balanced collections, and
the Bondareva–Shapley theorem~\citep{Bondareva1963, Shapley1967}
are not formalized.

\item \emph{Impossibility theorems in social choice.}  Arrow,
Gibbard–Satterthwaite, and Sen's liberal paradox are formalizable
in this framework but are not in the current library.

\item \emph{Bayesian persuasion and information design.}
Kamenica–Gentzkow~\citep{KamenicaGentzkow2011} and successors are
out of scope.

\item \emph{Mean-field / population games.}  These require the
continuum and are out of scope.
\end{itemize}

These are not bugs; they are the boundary of the library's
contract.
